\batchmode
\makeatletter
\def\input@path{{C:/Users/cash/Documents/}}
\makeatother
\documentclass[spanish]{report}
\usepackage{fontspec}
\setcounter{secnumdepth}{3}
\setcounter{tocdepth}{3}

\makeatletter

%%%%%%%%%%%%%%%%%%%%%%%%%%%%%% LyX specific LaTeX commands.
\providecommand\textquotedblplain{%
  \bgroup\addfontfeatures{RawFeature=-tlig}\char34\egroup}

\makeatother

\usepackage{babel}
\begin{document}
\title{Plataforma Evalis: Sistema de Gestión de Evaluaciones y Control Docente}
\author{Piero Stephano Rosario Sánchez}
\maketitle

\chapter{Introducción y Objetivos}
\begin{verse}
En el entorno educativo actual, la digitalización de los procesos
administrativos y académicos se ha vuelto una necesidad fundamental
para optimizar el tiempo del personal docente y directivo. El registro
de calificaciones, el control del profesorado y la generación de actas
oficiales suelen ser tareas complejas que, si se realizan de forma
manual o descentralizada, son propensas a errores y retrasos.
\end{verse}

\chapter{1. Base de datos}

\section{Diseño de la base de datos relacional}

Para garantizar la integridad, consistencia y escalabilidad de la
Plataforma Evalis, se ha diseñado una base de datos relacional robusta
utilizando el sistema de gestión de bases de datos MySQL. La estructura
de almacenamiento se compone de un conjunto de tablas interconectadas
mediante claves primarias y foráneas, permitiendo reflejar fielmente
la jerarquía y el flujo operativo de un centro educativo.

A continuación, se describen los módulos y tablas principales que
conforman el esquema del sistema:

\subsection{Gestión de Usuarios y Roles de Personal}

El núcleo del control de acceso y de la información biográfica se
compone de tablas genéricas especializadas a través de relaciones
heredadas o de cardinalidad uno a uno:
\begin{description}
\item [{persona:}] Tabla base que centraliza la información común de todos
los individuos del centro (DNI, nombre, apellidos, datos de contacto).
Evita la duplicidad de datos en el sistema.
\item [{estudiants}] / \textbf{professor} / \textbf{administradors} / \textbf{directiva}:
Tablas especializadas que se vinculan directamente con la tabla persona
mediante el uso del DNI como clave. Almacenan información exclusiva
de cada rol, como el NIA en el caso de los estudiantes o las competencias
contractuales en el de los profesores.
\item [{usuari}] / \textbf{login\_logs}: Entidades encargadas de gestionar
las credenciales de autenticación cifradas, el control de las sesiones
activas (sessions) y el registro de auditoría de los accesos al sistema.
\end{description}

\subsection{Estructura académica y plan de estudios}

Para la organización de las clases y los contenidos didácticos que
se imparten, se han implementado las siguientes relaciones:
\begin{description}
\item [{grup\_classe:}] Define los diferentes grupos o aulas del centro
(por ejemplo, 1r BATX A).
\item [{assignatura}] \textbf{/ modul / unitat:} Representan la jerarquía
del plan docente. Un módulo se compone de diversas asignaturas, las
cuales a su vez se subdividen en unidades didácticas configurables.
\item [{grupclasse\_assignatura}] / professor\_grup\_classe: Tablas intermedias
que resuelven las relaciones de muchos a muchos. Permiten asociar
qué asignaturas se imparten en qué grupos y qué profesores tienen
asignados dichos grupos.
\end{description}

\subsection{Evaluación, Asistencia e Incidencias}

El registro diario del rendimiento y el comportamiento del alumnado
se consolida a través de un bloque de tablas operativas:
\begin{description}
\item [{notes:}] Es la tabla central del módulo de evaluación. Almacena
las calificaciones asociando el identificador de la unidad didáctica
(uf\_id), el NIA del estudiante y el valor numérico de la nota.
\item [{estado\_evaluación:}] Controla de manera global si el sistema se
encuentra en un periodo abierto para la inserción de notas o si está
bloqueado.
\item [{gestio\_absencies}] / \textbf{asistencia\_per\_hora: }Registran
las faltas de asistencia y los retrasos cometidos por los alumnos
en franjas horarias específicas.
\item [{observacions\_incidencies:}] Espacio dedicado para que el equipo
docente anote amonestaciones o comentarios de seguimiento pedagógico.
\end{description}

\section{Dificultades Encontradas y Soluciones Aplicadas}

El desarrollo del modelo de datos de la Plataforma Evalis supuso un
reto organizativo y técnico considerable. En primer lugar, cabe destacar
que la planificación inicial del proyecto estaba diseñada para ser
ejecutada por equipos de dos integrantes. Al tener que asumir el desarrollo
del sistema de manera completamente individual, la carga de trabajo
en el diseño y la normalización de la base de datos se duplicó, obligando
a priorizar la eficiencia y a resolver de forma autónoma diversos
problemas técnicos complejos.

A continuación, se detallan los principales inconvenientes surgidos
durante la elaboración de la base de datos y cómo se solventaron:
\begin{description}
\item [{Problema}] con el diseño de claves primarias compuestas: Al modelar
la tabla notes, inicialmente se planteó usar un ID autonumérico único
como clave primaria. Sin embargo, esto permitía que el sistema pudiera
registrar por error más de una nota para un mismo alumno en la misma
unidad formativa.
\item [{Solución:}] Se reestructuró la tabla aplicando una clave primaria
compuesta por la combinación de nia (alumno) y uf\_id (unidad). De
este modo, el propio motor de MySQL impide la duplicidad de registros
a nivel físico y garantiza la consistencia de los datos históricos.
\item [{Sincronización}] de datos heredados (Especialización de roles):
El diseño requería que los profesores, estudiantes y administradores
compartieran datos comunes (como nombre, apellido o DNI) pero mantuvieran
campos exclusivos de su rol. Al principio, la separación de las tablas
(persona, estudiants, professor) provocaba problemas de sincronización
si se borraba o modificaba un usuario.
\item [{Solución:}] Se implementó una arquitectura de herencia en la base
de datos, donde el DNI actúa como Clave Primaria en persona y, a su
vez, como Clave Foránea (FK) con restricción de borrado y actualización
en cascada (ON DELETE CASCADE / ON UPDATE CASCADE) en las tablas hijas.
Esto asegura que cualquier cambio en los datos maestros se replique
instantáneamente en todo el sistema.
\item [{Inconsistencias}] al recuperar datos debido a valores nulos (NULL):
Al realizar las primeras pruebas de carga para el boletín de calificaciones,
la consulta fallaba o ignoraba a los alumnos que todavía no tenían
ninguna nota registrada en la tabla notes, ya que se utilizaba un
enlace de tipo INNER JOIN.
\item [{Solución:}] Se sustituyó la lógica de las consultas por un LEFT
JOIN combinado con la función IsDBNull en el entorno de desarrollo
(VB.NET). Con esto se logró que la aplicación liste siempre a la totalidad
de los alumnos matriculados en el grupo, mostrando la calificación
en blanco si aún no ha sido evaluado, en lugar de romper el flujo
de la pantalla.
\end{description}

\chapter{Desarrollo de la Aplicación Móvil}

\section{Enfoque Centrado en la Experiencia de Usuario (UX)}

A diferencia del módulo de escritorio de la Plataforma Evalis (diseñado
para la gestión operativa y carga masiva de datos por parte del personal
docente), la aplicación móvil se concibió única y exclusivamente como
un portal de servicios orientado al estudiante. El objetivo principal
de este desarrollo es empoderar al alumno, ofreciéndole una interfaz
moderna, intuitiva y accesible que centralice toda su vida académica
en la palma de su mano.

Para lograr una experiencia de usuario (UX) óptima y adaptada a los
estándares actuales, la aplicación implementa las siguientes características
de diseño y accesibilidad:
\begin{description}
\item [{Personalización}] de interfaz (Tema Oscuro / Negro): Para mejorar
el confort visual en entornos de baja luminosidad y optimizar el consumo
de batería en pantallas de tecnología OLED/AMOLED, se desarrolló un
sistema de gestión de temas que permite al usuario alternar entre
el modo claro tradicional, un tema oscuro y un modo negro puro según
sus preferencias.
\item [{Navegación}] Fluida y Adaptable: Se sustituyeron las estructuras
rígidas y los grandes contenedores de escritorio por un diseño móvil
nativo basado en paneles limpios y menús deslizantes, garantizando
que la consulta de datos sea rápida, intuitiva y libre de fricciones.
\end{description}

\section{Módulos y Funcionalidades de la App Móvil}

Al estar la aplicación móvil enfocada por completo en la experiencia
del alumno, todas sus funciones se estructuran para facilitarle el
acceso a la información y agilizar la interacción con el centro. Los
módulos principales implementados son los siguientes:
\begin{description}
\item [{Panel}] de Calificaciones y Notificaciones en Tiempo Real: El alumno
puede consultar sus notas actualizadas al instante. Además, el sistema
integra un servicio de alertas que envía una notificación push al
dispositivo móvil en el momento exacto en que un profesor publica
o modifica una calificación en la base de datos.
\item [{Módulo}] de Profesorado y Búsqueda en Tiempo Real: Para facilitar
la localización y el contacto con los docentes, la aplicación incluye
una pantalla que lista a los profesores vinculados al alumno, organizándolos
de forma clara según la asignatura específica que le imparten. Además,
se ha implementado una barra de búsqueda inteligente con filtrado
en tiempo real, lo que permite encontrar a cualquier profesor al instante
conforme se va tecleando su nombre.
\item [{Estadísticas}] Académicas y Rendimiento: El módulo de analítica
calcula y muestra de forma gráfica la nota media del estudiante por
asignatura, permitiéndole realizar un seguimiento visual de su progreso
y rendimiento a lo largo del curso.
\item [{Control}] de Prácticas en Empresa (FCT): Espacio dedicado al seguimiento
del módulo de Formación en Centros de Trabajo, donde el alumno puede
monitorizar de forma exacta el cómputo acumulado de sus horas de prácticas
realizadas en la empresa colaboradora.
\item [{Expediente}] e Información Académica: Permite el acceso al historial
de estudios finalizados de cursos anteriores, así como la consulta
de datos esenciales del perfil del alumno y el aula asignada.
\item [{Generación}] y Descarga de Boletín en PDF: Mediante la integración
de un motor ligero de renderizado de documentos, el estudiante puede
exportar e imprimir su boletín oficial de calificaciones directamente
en formato PDF desde el propio terminal, sin perder el formato institucional.
\end{description}

\section{Dificultades Técnicas en el Desarrollo Individual}

Desarrollar una aplicación con tal densidad de funcionalidades de
manera completamente individual (asumiendo en solitario el trabajo
originalmente planificado para un equipo de dos personas) supuso afrontar
retos técnicos complejos, especialmente en la gestión y optimización
de los datos:
\begin{description}
\item [{Optimización}] y rendimiento en la carga del expediente completo:
Al abrir el perfil o el expediente, la aplicación móvil debe realizar
múltiples consultas a la base de datos MySQL para traer a la vez las
notas actuales, las medias calculadas, el estado de las FCT y los
datos del aula. Hacer esto con peticiones separadas colapsaba la conexión
y hacía que la app tardara varios segundos en cargar la pantalla.
\item [{Solución:}] Se diseñaron consultas optimizadas utilizando vistas
en la base de datos y uniones avanzadas (LEFT JOIN). De este modo,
el servidor procesa toda la información en una sola petición estructurada,
consiguiendo que el expediente y las estadísticas del alumno se muestren
de forma instantánea al abrir la app.
\item [{Sincronización}] del sistema de alertas sin saturar el servidor:
Conseguir que el móvil se entere de una nueva nota sin saturar la
base de datos con peticiones constantes (polling) desde el dispositivo
del alumno era inviable para el rendimiento del sistema.
\item [{Solución:}] Se diseñó una lógica eficiente basada en eventos. La
aplicación móvil comprueba de forma inteligente los cambios de fecha
y estado en las tablas de la base de datos para disparar la notificación
local únicamente cuando se detecta una inserción real en la tabla
notes que afecte al NIA del alumno autenticado.
\end{description}

\chapter*{Desarrollo de la Aplicación de Escritorio}

\section{Enfoque Arquitectónico y Modular}

La aplicación de escritorio de la Plataforma Evalis constituye el
núcleo administrativo y operativo del sistema. A diferencia de la
versión móvil (orientada exclusivamente a la consulta por parte del
alumnado), el entorno de escritorio está diseñado específicamente
para el personal docente, directivo y de administración. Su propósito
principal es ofrecer una herramienta de alto rendimiento que simplifique
la carga masiva de datos, el control de las evaluaciones y la gestión
global del centro.

Para garantizar la mantenibilidad del código y permitir futuras expansiones
sin necesidad de reescribir la aplicación, se ha optado por un desarrollo
modular en Visual Basic .NET utilizando Windows Forms. La interfaz
se articula mediante controles de usuario personalizados (UserControls)
que se cargan de forma dinámica en un contenedor principal, evitando
la apertura caótica de múltiples ventanas y unificando la experiencia
de usuario.

\section{Gestión de Roles y Control de Acceso}

Para garantizar la seguridad de la información académica y evitar
alteraciones no autorizadas en las calificaciones, la plataforma implementa
un sistema de control de acceso basado en roles. Dependiendo del tipo
de cuenta con la que el usuario inicie sesión en la aplicación, el
sistema habilitará o bloqueará dinámicamente las diferentes pestañas
y funciones de la interfaz.

A continuación, se describen las responsabilidades, funcionalidades
y privilegios asignados a cada perfil operativo de la versión de escritorio:

\subsection{Rol: Profesor}

Es el perfil operativo básico dentro del entorno docente. Sus funciones
están estrictamente limitadas a la gestión y evaluación de las asignaturas
que tiene asignadas en su horario lectivo:
\begin{description}
\item [{Volcado}] de Calificaciones mediante DataGridView: Permite introducir
y almacenar las notas de los alumnos matriculados en sus grupos de
clase de forma masiva a través de una tabla dinámica interactiva (dgvNotas).
\item [{Consulta}] del Claustro Docente: Dispone de una vista exclusiva
para listar y consultar a todos los profesores del centro con sus
respectivas especialidades.
\item [{Restricción}] Temporal de Edición: El profesor solo tiene permitido
insertar o modificar calificaciones cuando el periodo de evaluación
se encuentra en estado \textquotedbl ABIERTO\textquotedbl . Si el
periodo está cerrado, la interfaz bloquea automáticamente la edición
del DataGridView, impidiendo cualquier cambio físico en la tabla notes.
\end{description}

\subsection{Rol: Jefe de Estudios}

Este perfil dispone de permisos administrativos intermedios enfocados
en la supervisión académica del centro educativo y la gestión de los
plazos temporales:
\begin{description}
\item [{Control}] del Periodo Evaluativo: Es el encargado directo de modificar
la tabla estado\_evaluacion, abriendo o cerrando los plazos de entrega
de notas para todo el claustro docente con un solo clic.
\item [{Exportación}] de Actas Oficiales: Cuenta con la capacidad de generar
y exportar documentos con las calificaciones consolidadas de todos
los alumnos pertenecientes a cualquier clase del centro para su posterior
tramitación administrativa.
\item [{Privilegios}] Docentes Expandidos: Reúne las mismas capacidades
de inserción y edición de notas que un profesor (a través de la misma
funcionalidad de cuadrícula de datos), pero con la ventaja jerárquica
de poder editar calificaciones o supervisar actas bloqueadas si la
situación lo requiere.
\end{description}

\subsection{Rol: Administrador (Admin)}

Es el rol con mayor jerarquía dentro del ecosistema de la plataforma.
Posee control total sobre el software, las estructuras maestras y
las herramientas de auditoría de seguridad:
\begin{description}
\item [{Supervisión}] Total (Full Access): Hereda y puede ejecutar de forma
nativa la totalidad de las funciones de los roles de Profesor y Jefe
de Estudios de manera conjunta y sin ningún tipo de restricción.
\item [{Interoperabilidad}] de Datos Docentes (JSON): Dispone de una herramienta
exclusiva para realizar extracciones masivas de la tabla professor
y exportar la información completa del claustro en formato JSON para
integraciones externas.
\item [{Auditoría}] de Seguridad y Accesos (XML): Para garantizar la seguridad
del sistema frente a accesos indebidos o registrar acciones críticas,
el administrador puede realizar consultas avanzadas sobre la tabla
login\_logs y exportar el historial completo de conexiones de cualquier
persona del centro en un archivo estructurado XML.
\end{description}

\section{Dificultades Técnicas en el Desarrollo Individual}

El principal y más crítico desafío durante el desarrollo de la versión
de escritorio de la Plataforma Evalis fue el factor tiempo. Al tener
que asumir de forma estrictamente individual un proyecto cuya envergadura,
alcance y diseño estaban planificados originalmente para ser ejecutados
por un equipo de dos integrantes, el cronograma de trabajo se vio
severamente ajustado.

Esta limitación temporal obligó a aplicar una estrategia de desarrollo
basada en los siguientes puntos:
\begin{description}
\item [{Sobrecarga}] en la fase de desarrollo y maquetación: Compaginar
en solitario el diseño de la interfaz gráfica en Windows Forms, la
programación de la lógica en VB.NET y la estructuración de los roles
de acceso duplicó las horas de dedicación necesarias en comparación
con un entorno de trabajo colaborativo.
\item [{Solución}] y Mitigación: Se aplicó una estricta priorización de
funcionalidades (enfoque MVP o Producto Mínimo Viable), asegurando
que los módulos críticos ---como el volcado de notas mediante el
DataGridView y el control de accesos por rol--- estuvieran completamente
operativos y pulidos antes de avanzar con los extras del sistema.
Se optimizaron las jornadas de desarrollo mediante una planificación
modular para garantizar la entrega puntual del proyecto final dentro
del plazo establecido.
\end{description}

\chapter{Arquitectura de Comunicación y Seguridad (API PHP)}

La aplicación móvil no realiza conexiones directas a la base de datos
centralizada, protegiendo así la integridad del servidor MySQL. En
su lugar, se ha diseñado e implementado una arquitectura cliente-servidor
basada en una API de servicios web en PHP.

El flujo de comunicación y el intercambio de información sigue un
protocolo seguro estructurado en tres niveles:
\begin{enumerate}
\item Peticiones por URL: La aplicación móvil solicita información (como
el expediente, las notas o el perfil) realizando peticiones HTTP a
URLs específicas que apuntan a los archivos PHP alojados en el servidor.
\item Validación de Identidad mediante Tokens: Para evitar el robo de información
o accesos no autorizados, cada petición debe incluir un Token de Seguridad
único generado tras el inicio de sesión. Los archivos PHP interceptan
este token y comprueban su validez antes de procesar cualquier orden.
Si el token es inválido o ha expirado, el servidor deniega el acceso
automáticamente.
\item Procesamiento y Respuesta: Una vez superada la barrera de seguridad
del token, el script PHP interactúa de forma segura con MySQL, extrae
los datos solicitados y los envía de vuelta al dispositivo móvil de
manera estructurada.
\end{enumerate}

\section{Dificultades Técnicas en el Desarrollo Individual}

Desarrollar un ecosistema con este nivel de robustez y seguridad de
manera completamente individual (asumiendo en solitario el trabajo
de dos personas) supuso afrontar retos técnicos complejos en la sincronización
de las sesiones:
\begin{description}
\item [{Gestión}] y persistencia del Token de Seguridad en el dispositivo:
Al implementar la seguridad por tokens, un problema recurrente era
que el alumno perdía el acceso o tenía que loguearse continuamente
cada vez que la app móvil cambiaba de pantalla o se minimizaba, debido
a la volatilidad de la memoria del terminal.
\item [{Solución:}] Se programó un sistema de almacenamiento seguro y local
en la aplicación móvil que retiene el token encriptado de forma persistente
mientras la sesión esté activa. La app adjunta este token de forma
automática en las cabeceras de las URLs de consulta hacia los archivos
PHP, garantizando una navegación fluida, segura y transparente para
el usuario.
\item [{Sincronización}] del sistema de alertas sin saturar el servidor:
Conseguir que el móvil se entere de una nueva nota sin saturar la
base de datos con peticiones constantes (polling) desde el dispositivo
del alumno era inviable para el rendimiento del sistema.
\item [{Solución:}] Se diseñó una lógica eficiente basada en eventos. La
aplicación móvil comprueba de forma inteligente los cambios de fecha
y estado a través de la API en PHP para disparar la notificación local
únicamente cuando se detecta una inserción real en la tabla notes
que afecte al NIA del alumno autenticado.
\end{description}

\section{Estrategia de Mitigación y Priorización del Alcance}

Como solución general ante la falta de tiempo derivada de asumir el
proyecto de manera estrictamente individual (cuando estaba planificado
para dos personas), se adoptó una metodología de desarrollo ágil orientada
a la priorización del alcance. En lugar de intentar abarcar todas
las funcionalidades de manera superficial, la estrategia se centró
en asegurar la máxima robustez y un acabado profesional en el núcleo
del ecosistema de la Plataforma Evalis.

Para gestionar de forma eficiente el calendario de entregas, se clasificaron
los objetivos del desarrollo bajo los siguientes criterios de prioridad:
\begin{description}
\item [{Elementos}] de Prioridad Crítica (Desarrollo Completo): Se priorizó
el diseño físico de la base de datos centralizada en MySQL, la arquitectura
de seguridad basada en Tokens a través de la API en PHP para la app
móvil, y el control estricto de accesos por roles en el entorno de
escritorio de VB.NET. De igual forma, se aseguró el correcto funcionamiento
del volcado masivo de calificaciones en el DataGridView y la descarga
del boletín en PDF desde el terminal móvil, al ser las herramientas
esenciales para el flujo de trabajo del centro educativo.
\item [{Elementos}] de Prioridad Secundaria (Postergados para Líneas Futuras):
Aquellas funcionalidades de carácter cosmético, los sistemas de mensajería
interna no críticos o las optimizaciones avanzadas de la interfaz
visual que requerían una alta inversión de tiempo y no aportaban valor
inmediato a la operatividad del sistema, se delegaron deliberadamente
para futuras actualizaciones del software.
\end{description}
Esta toma de decisiones estratégica permitió mitigar eficazmente el
impacto del tiempo en contra, garantizando que el producto final entregado
sea un software real, estable, completamente seguro y funcional en
todas sus plataformas, cumpliendo rigurosamente con los estándares
de calidad exigidos.

\bibliographystyle{plain}

\end{document}
